\section{Discussion}

\subsection{What closed-loop reasoning contributes}

The central finding of this study is that closing the feedback loop between LLM planning and MLFF relaxation yields three key advantages: reliability (including cross-backend convergence), self-reflection, and interpretability. Previously, each property had shown powers in other frameworks, but no existing adsorption search system successfully integrated all three simultaneously. For instance, active learning enhances reliability~\cite{settles2012active}, MLFFs decrease per-evaluation costs~\cite{lan2023adsorbml}, and diagnostic workflows offer interpretability.

\textbf{Reliability.} Across both benchmarks and all four LLM backends, \textit{AdsMind Full} achieves near-perfect success rates; the few observed failures stem from genuine chemical dissociation events rather than framework errors. DFT validation confirms sign-correctness across all six tested systems. Furthermore, iterative feedback drives heterogeneous LLM backends toward consistent physical solutions (achieving over a 2$\times$ variance reduction on CMU20 and over 1.7$\times$ on OCD-GMAE62). This ensures that conclusions remain independent of the specific LLM employed, which is an important feature for LLM-based architecture, which relies on hosted, black-box models subject to unannounced updates.

% reflection in planner
\textbf{Efficiency.} \textit{AdsMind Full} requires approximately 4 relaxations per case, representing a 5--25$\times$ reduction in computational cost compared to the baselines evaluated in Section~\ref{sec:results_baselines}. This cost improvement fundamentally shifts the paradigm of viable search strategies: at $\sim$4 relaxations, closed-loop reasoning becomes computationally competitive with single-point evaluations. %Consequently, it can be seamlessly integrated into high-throughput screening pipelines where brute-force enumeration remains cost-prohibitive.

\textbf{Interpretability.} The closed-loop architecture yields physically grounded artifacts, such as slip events, FORBID constraints, and site-mismatch fingerprints, that are inaccessible to open-loop systems. By stratifying chemical slip rates across surface families, we establish a quantitative proxy for evaluating an LLM's geometric reasoning regarding binding sites, offering a valuable complement to text-based benchmarks like ChemBench~\cite{mirza2024chembench}. Unlike conventional evaluation suites that primarily test declarative chemical knowledge, slip tracking reveals whether an LLM's geometric intuition withstands the rigors of physical relaxation---a distinct capability strictly necessary for autonomous scientific workflows.

\subsection{Limitations}

\textbf{Energy depth.} For complex adsorbates (e.g., CH$_2$CH$_2$OH, OCHCH$_3$), AdsMind Full occasionally yields higher final energies than broad-sampling baselines. While multi-seed execution (5 seeds for cases 16--18) recovers a portion of this depth (gains of $+0.038$ to $+0.524$~eV), a residual gap of $\sim$1~eV remains relative to Adsorb-Agent. This discrepancy is a consequence of architectural design rather than a functional defect: iterative single-trajectory refinement is optimized for local basin convergence and reliability, rather than exhaustive global exploration.

\textbf{FORBID over-pruning.} The use of explicit negative constraints can occasionally lead to over-pruning, particularly when an initial placement is fortuitously optimal. This is evidenced by case 008 of OCD-GMAE62, where the \textit{Full} variant underperforms relative to other versions. Future work should investigate adaptive constraints that dynamically relax after detecting convergence to mitigate such over-exclusion.

\textbf{Force-field scope.} Our evaluation relies on MACE-MP-0 small as the primary physical backend. Sensitivity checks with the "large" variant on CMU20 reveal a mean absolute shift of 1.255~eV, highlighting that absolute rankings may vary with the chosen MLFF. While AdsMind correctly identifies exothermic binding across all DFT-validated systems (Section~\ref{sec:dft_reference}), its accuracy is inherently bounded by the underlying MLFF's handling of dispersion, strong correlation, or charge-transfer effects. Extending AdsMind to these regimes requires higher-fidelity potentials but necessitates no structural changes to the agent architecture.

% limitation of FF
% \textbf{System Coverage.} The current implementation is validated for flat and stepped crystalline surfaces with single adsorbates. Although extensions to defects, nanoparticles, and co-adsorption are architecturally straightforward, they require further rigorous validation against reference datasets such as QUASAR~\cite{yang2026quasar}.

% \textbf{Statistical heterogeneity.} While Friedman tests confirm overall significance ($p < 0.01$), many per-backend Wilcoxon tests do not survive Benjamini-Hochberg (BH) correction. Given this heterogeneity, we prioritize effect sizes and bootstrap intervals over binary significance thresholds to provide a more robust quantification of framework performance.

\textbf{LLM reproducibility.} Our Tier 2 stability audit ($N=3$) reveals that while 72.9\% of runs are perfectly reproducible (range $<$0.001~eV), a stochastic tail (~21\%) exists where LLM non-determinism drives multi-basin exploration. This confirms that single-run energies should be interpreted as representative samples rather than unique global solutions. We mitigate this through a multi-backend ensemble design and the use of frozen model identifiers.

% evolution of cross session knowledge

\section{Conclusion}

We introduce AdsMind, a closed-loop multi-agent framework that couples LLM-driven adsorption hypothesis generation with iterative MLFF relaxation feedback. This approach overcomes the fundamental limitation of open-loop agents: the inability to detect, diagnose, and recover from reasoning errors using physical evidence.
Evaluated across 20 (CMU20) and 62 (OCD-GMAE62) absorption processes utilizing four diverse LLM backends, AdsMind simultaneously delivers reliability, efficiency, and interpretability. This is a combination unachieved by prior open-loop adsorption search systems. By deliberately occupying the reliability-first, low-computational-cost quadrant of the search space, AdsMind serves as a powerful complement to depth-maximizing enumeration strategies.
Future work should focus on integrating AdsMind's closed-loop error recovery with broader initial enumeration, alongside extending validation to more complex environments such as co-adsorption, solvation, and defected surfaces. More broadly, the Chemical slip diagnostic, which quantifies the mismatch between LLM-planned and physically realized binding sites, establishes a robust template for grounding language-model reasoning in physical reality, carrying significant implications for autonomous scientific discovery across computational chemistry.